% \enabletrackers[sql*]
% \enabletrackers[*lib*]

% E:\tex-context\tex\texmf-win64\bin\lib\luametatex\sqlite\sqlite3.dll

% with sqlite we have only single access (unless we check for locks and
% wait a while .. todo)
%
% with a mysql server one can access in parallel and make a
% high performance parallel processor: fetch what is needed

% \registerctxluafile{libs-imp-sqlite}{autosuffix}

\usemodule[sql][method=sqlite]

\startluacode

    local sql = require("util-sql")

    document.sqlconverter = sql.makeconverter {
        { name = "index", type = "number" },
        { name = "x",     type = "number" },
        { name = "y",     type = "number" },
    }

 -- sql.setmethod("sqlite") -- mysql

    document.sqlpresets = {
        host      = "localhost",
        username  = "root",
        password  = "secret",
        database  = "test",
        id        = "test"
    }

    local sqlite_template = [[
        CREATE TABLE IF NOT EXISTS `test` (
            `index`  INTEGER PRIMARY KEY,
            `x`      INTEGER,
            `y`      INTEGER
        )
    ]]

    local results, rest = utilities.sql.execute {
        presets   = document.sqlpresets,
        template  = sqlite_template,
    }

 -- inspect(results)
 -- inspect(rest)

\stopluacode

\startluacode
    local template = [[
        INSERT INTO `test` (
            `index`,
            `x`,
            `y`
        ) VALUES (
            %[index]%,
            %[x]%,
            %[y]%
        ) ;
    ]]

    local function add(t)
        local results = utilities.sql.execute {
            presets   = document.sqlpresets,
            template  = template,
         -- variables = {
         --     index = t.index,
         --     x     = t.x,
         --     y     = t.y,
         -- },
            variables = t,
        }
        if results then
--            inspect(results)
        end
    end

--  add { index = 1, x = 10, y = 30 }
--  add { index = 2, x = 20, y = 20 }
--  add { index = 3, x = 30, y = 10 }

    function mp.AddPoint(index, x, y)
        add { index = index, x = x, y = y }
    end

\stopluacode

\startluacode
    local template = [[
        UPDATE `test`
        SET
            `x` = %[x]%,
            `y` = %[y]%
        WHERE
            `index` = %[index]%
        ;
    ]]

    local function update(t)
        local results, x, y = utilities.sql.execute {
            presets   = document.sqlpresets,
            template  = template,
            variables = t,
        }
        if results then
--            inspect(results)
        end
    end

--  update { index = 2, x = 22.22, y = 2.2 }

    function mp.UpdatePoint(index, x, y)
        update { index = index, x = x, y = y }
    end
\stopluacode

\startluacode

    local template = [[
        SELECT * from `test`
        WHERE
            `index` = %[index]%
        ;
    ]]

    local function fetch(index)
        local results = utilities.sql.execute {
            presets   = document.sqlpresets,
            template  = template,
            variables = { index = index },
            converter = document.sqlconverter,
        }
        local point = results and results[1]
        if point then
            return { point.x or 0, point.y or 0 }
        else
            return { 0, 0 }
        end
    end

    function mp.FetchPoint(index)
        mp.inject(fetch(index))
    end

\stopluacode

\usemodule[article-basic]

\startluacode

    metapost.registerdirect("FetchPoint", function()
        return mp.FetchPoint(
            mp.scan.integer()
        )
    end)
    metapost.registerdirect("AddPoint", function()
        return mp.AddPoint(
            mp.scan.integer(),
            mp.scan.integer(),
            mp.scan.integer()
        )
    end)
    metapost.registerdirect("UpdatePoint", function()
        return mp.UpdatePoint(
            mp.scan.integer(),
            mp.scan.integer(),
            mp.scan.integer()
        )
    end)

\stopluacode

\startMPdefinitions
    newscriptindex mfid_FetchPoint  ; mfid_FetchPoint  := scriptindex "FetchPoint"  ;
    newscriptindex mfid_AddPoint    ; mfid_AddPoint    := scriptindex "AddPoint"    ;
    newscriptindex mfid_UpdatePoint ; mfid_UpdatePoint := scriptindex "UpdatePoint" ;

    vardef FetchPoint(expr n) =
        runscript mfid_FetchPoint n
    enddef ;
    vardef AddPoint(expr n, x, y) =
        runscript mfid_AddPoint n x y
    enddef ;
    vardef UpdatePoint(expr n, x, y) =
        runscript mfid_UpdatePoint n x y
    enddef ;
\stopMPdefinitions

\starttext

\startMPpage
    pair ThePosition[];
    ThePosition[1] := lua.mp.FetchPoint(10) ;
    ThePosition[2] := lua.mp.FetchPoint(15) ;
    draw (ThePosition[1] -- ThePosition[2]) scaled 20 ;
    % first time
lua.mp.AddPoint(10,2,3) ;
lua.mp.AddPoint(15,5,6) ;
    % after that
    lua.mp.UpdatePoint(10,2,3.3) ;
    lua.mp.UpdatePoint(15,5,6) ;
\stopMPpage

% \startMPpage
%     pair ThePosition[];
%     ThePosition[1] := FetchPoint(10) ;
%     ThePosition[2] := FetchPoint(15) ;
%     draw (ThePosition[1] -- ThePosition[2]) scaled 20 ;
% %     % first time
% % AddPoint(10,2,3) ;
% % AddPoint(15,5,6) ;
%     UpdatePoint(10,2,3.3) ;
%     UpdatePoint(15,5,6) ;
% %     % after that
% \stopMPpage

\startluacode

    context.starttitle { title = "Points" }

    local results = utilities.sql.execute {
        presets   = document.sqlpresets,
        template  = "select * from `test` ;",
        variables = { },
        converter = document.sqlconverter,
    }

 -- inspect(results)

    context.starttabulate { "|c|r|r|" }
    context.HL()
    context.NC() context.bold("index")
    context.NC() context.bold("x")
    context.NC() context.bold("y")
    context.NC() context.NR()
    context.HL()
    for i=1,#results do
        local result = results[i]
        context.NC() context(result.index)
        context.NC() context(result.x)
        context.NC() context(result.y)
        context.NC() context.NR()
    end
    context.HL()
    context.stoptabulate()

    context.stoptitle()

\stopluacode

\stoptext

% just for the fun of it
%
% store (database,listname,value)
%
% numeric -> native
% boolean -> native
% pair    -> native
% pair    -> native
% path    -> serialized lua
%
% bump performance with by using indexed runscript
